\documentclass{subfiles}
\begin{document}
    Sometimes referred to as \gls{eca}, Lenstra algorithm is a way to test whether a number \(n\)
        is prime, and eventually find a non-trivial divisor.
        The idea is the following, we select a curve \(\C\) and a point \(P \in \C\),
        and we compute all multiples of \(P\) up to a theshold \(B\), by reducing modulo \(n\).
        If it happens that at a certain point we can't compute the inverse of the divisor,
        then we'll have found a non-trivial divisor of \(n\).
        More precisely, if \(v\) is the our non-invertible element, 
        then \(\gcd(n, v)\) is the non-trivial divisor of \(n\).

    It can happen that no non-invertible element is found,
        in that case we should either choose a different point,
        a different curve or both.
\end{document}

